\documentclass[12pt, mexican]{report}
\usepackage[utf8]{inputenc}
\usepackage[T1]{fontenc}
\usepackage{lmodern}
\usepackage{babel}
\usepackage{amsmath}
\usepackage{amsfonts}
\usepackage{amssymb}
\usepackage{geometry}
\usepackage{titlesec}
\usepackage{graphicx}
\usepackage{hyperref}
\usepackage{xcolor}
\usepackage{listings}

\lstset{
    language=Python,
    basicstyle=\ttfamily\footnotesize,
    keywordstyle=\color{blue},
    stringstyle=\color{teal},
    commentstyle=\color{green!50!black},
    numbers=left,
    numberstyle=\tiny,
    stepnumber=1,
    numbersep=5pt,
    backgroundcolor=\color{white},
    showspaces=false,
    showstringspaces=false,
    showtabs=false,
    frame=single,
    tabsize=4,
    captionpos=b,
    breaklines=true,
    breakatwhitespace=false,
    extendedchars=true,
    literate={á}{{\'a}}1 {é}{{\'e}}1 {í}{{\'i}}1 {ó}{{\'o}}1 {ú}{{\'u}}1 {Á}{{\'A}}1 {É}{{\'E}}1 {Í}{{\'I}}1 {Ó}{{\'O}}1 {Ú}{{\'U}}1 {ñ}{{\~n}}1 {Ñ}{{\~N}}1
}

\geometry{a4paper, margin=2.5cm}

% Configuración de títulos para que parezcan secciones de un reporte
\titleformat{\chapter}[display]
  {\normalfont\bfseries}{}{0pt}{\Huge}
\titlespacing*{\chapter}{0pt}{-20pt}{20pt}

\title{
    \textbf{Universidad Nacional Autónoma de México} \\
    \Large Facultad de Ciencias \\
    \vspace{1cm}
    \large Análisis Matemático Aplicado \\
    \vspace{0.5cm}
    \textbf{Tarea 2: Optimización No Lineal y Estimación de Parámetros usando el Algoritmo Adam}
}
\author{Karla Romina Juárez Torres}
\date{4 de mayo del 2026}

\begin{document}

\maketitle

\chapter{Introducción}
Implementación y análisis del algoritmo de optimización Adam para resolver un problema de regresión no lineal basado en el modelo:
\begin{equation}
    f(x;m)=a+\frac{\sin(2\pi bx)}{x^{c}}
\end{equation}

\section{Inciso A: Derivación Analítica del Gradiente (Regla de la Cadena)}

Para optimizar $m=[a,b,c]^{T}$ con Adam, necesitamos el gradiente de la función de pérdida MSE:
\begin{equation}
    L(m)=\frac{1}{2N}\sum_{i=1}^{N}(f(x_{i};m)-y_{i})^{2}
\end{equation}

\subsection{Derivadas de \texorpdfstring{$f(x;m)$}{f(x;m)}}

Calculamos las derivadas parciales respecto a cada parámetro:

\begin{enumerate}
    \item Respecto a $a$:
    \begin{equation}
        \frac{\partial f}{\partial a} = 1
    \end{equation}

    \item Respecto a $b$:
    \begin{equation}
        \frac{\partial f}{\partial b} = \frac{1}{x^{c}} \cdot \cos(2\pi bx) \cdot (2\pi x) = \frac{2\pi x \cos(2\pi bx)}{x^{c}}
    \end{equation}

    \item Respecto a $c$:
    \begin{equation}
        \frac{\partial f}{\partial c} = \sin(2\pi bx) \cdot \left( -x^{-c} \ln(x) \right) = -\frac{\sin(2\pi bx) \ln(x)}{x^{c}}
    \end{equation}
\end{enumerate}

\subsection{Gradiente \texorpdfstring{$\nabla L(m)$}{grad L(m)}}

Por regla de la cadena, el gradiente de $L(m)$ es:
\begin{equation}
    \frac{\partial L}{\partial m_{j}} = \frac{1}{N} \sum_{i=1}^{N} (f(x_{i};m) - y_{i}) \frac{\partial f(x_{i};m)}{\partial m_{j}}
\end{equation}

Sustituyendo las derivadas parciales, obtenemos el vector gradiente:
\begin{equation}
    \nabla L(m) = \begin{bmatrix}
        \frac{1}{N} \sum_{i=1}^{N} (f(x_{i};m) - y_{i}) \\
        \frac{1}{N} \sum_{i=1}^{N} (f(x_{i};m) - y_{i}) \left( \frac{2\pi x_{i} \cos(2\pi bx_{i})}{x_{i}^{c}} \right) \\
        \frac{1}{N} \sum_{i=1}^{N} (f(x_{i};m) - y_{i}) \left( -\frac{\sin(2\pi bx_{i}) \ln(x_{i})}{x_{i}^{c}} \right)
    \end{bmatrix}
\end{equation}


\section{Inciso B: Generación de Datos (Ground Truth)}

Se generaron $200$ datos sintéticos equiespaciados en $x \in [0.5, 10]$ y posteriormente se evaluó $f(x; m)$ con parámetros reales $m_{real} = [2.0, 1.5, 0.5]^T$ para obtener $y_{limpio}$. 

Finalmente se usó ruido gaussiano ($\mu=0, \sigma=0.5$) para simular condiciones reales:
\begin{equation}
    y = y_{limpio} + \mathcal{N}(0, 0.5^2)
\end{equation}

Código base:

\lstinputlisting[language=Python, caption={Generacion\_datos.py}]{../Scripts/Generacion_datos.py}
\section{Inciso C: Implementación del Optimizador Adam}
\textit{[Espacio reservado para la explicación de la implementación en Python/NumPy].}


\section{Inciso D: Análisis Visual y Numérico}
\textit{[Espacio reservado para las gráficas de convergencia y ajuste de curva].}

\end{document}